\subsection{Summary}
The vertex and energy reconstruction of low--energy neutrino interactions in DUNE has been studied. The vertex resolution was evaluated using a point--like reconstruction approach that combines the TPC cluster vertex with the associated PDS flash. The achieved vertex resolution is between \SI{0.5}{\centi\meter} and \SI{3.5}{\centi\meter}, depending on the coordinate and the presence of background events. There are significant differences in fiducialization performance between the \ac{HD} and \ac{VD} geometries, with the \ac{HD} geometry achieving higher fiducialization efficiency due to its lower background rates. The energy reconstruction was performed using four different algorithms to sum the energy of reconstructed \ac{TPC} clusters. The performance of these algorithms was evaluated in terms of energy resolution, fiducialization efficiency, and the impact of flash--matching on the energy reconstruction. Contrary to the vertex resolution, the energy resolution does not show significant differences between the two geometries provided that one choses the appropiate reconstruction algorithm. The best energy resolution is achieved using the \textit{SolarEnergy} algorithm, which yields a resolution of approximately \num{20}\% at \SI{15}{\mega\electronvolt} and improves to about \num{10}\% at \SI{5}{\mega\electronvolt}. The flash--matching step has a minor impact on the energy resolution, indicating that other factors, such as clustering inefficiencies and intrinsic energy losses, dominate the resolution. Overall, the reconstruction algorithms demonstrate good performance for low--energy neutrino interactions in DUNE, with potential for further optimization and improvement.

Based on the amount of background activity in the \ac{VD} module, other reconstruction objectives such as channel tagging (out of the scope of this work) and particle identification are likely to be more challenging than in the \ac{HD} module. The case of a shielded \ac{VD} module, which has background rates comparable to the \ac{HD} module, is therefore highly encouraged and expected to yield similar low--energy reconstruction performance.